\documentclass[11pt,a4paper]{article}

% ---------- Core encoding / fonts ----------
\usepackage[utf8]{inputenc}
\usepackage[T1]{fontenc}
\usepackage{lmodern}
\usepackage{microtype}
\usepackage[strings]{underscore}   % literal underscores in text mode, no manual escaping

% ---------- Layout ----------
% Generous top/bottom margins so the header/footer sit clearly between
% the page-edge border (drawn 1cm in, see below) and the body text --
% no overlap in either direction.
\usepackage[a4paper,left=2.6cm,right=2.6cm,top=3.4cm,bottom=3.2cm,
            headheight=14pt,headsep=22pt,footskip=28pt]{geometry}
\usepackage{tikz}
\usetikzlibrary{calc}
\usepackage{eso-pic}
\usepackage{fancyhdr}
\usepackage{titlesec}
\usepackage{enumitem}
\usepackage{listings}
\usepackage{xcolor}
\usepackage{hyperref}
\usepackage{parskip}
\usepackage{seqsplit}
\usepackage{booktabs}
\usepackage{longtable}
\usepackage{array}

% ---------- Strictly black & white ----------
\hypersetup{
  colorlinks=true,
  linkcolor=black,
  citecolor=black,
  urlcolor=black,
  pdftitle={Kernel Borderlands -- Session Report, 2026-10-31/11-01},
  pdfauthor={Kernel Borderlands Team}
}

% ---------- Four-side page border on every page ----------
\newcommand{\pageborder}{%
  \begin{tikzpicture}[remember picture, overlay]
    \draw[black, line width=0.7pt]
      ($(current page.north west)+(1.0cm,-1.0cm)$) rectangle
      ($(current page.south east)+(-1.0cm,1.0cm)$);
  \end{tikzpicture}%
}
\AddToShipoutPictureBG{\pageborder}

% ---------- Headers / footers ----------
\pagestyle{fancy}
\fancyhf{}
\renewcommand{\headrulewidth}{0.4pt}
\renewcommand{\footrulewidth}{0pt}
\fancyhead[L]{\small\scshape Kernel Borderlands}
\fancyhead[R]{\small\scshape Session Report}
\fancyfoot[C]{\small\thepage}

\fancypagestyle{plain}{%
  \fancyhf{}
  \renewcommand{\headrulewidth}{0pt}
  \fancyfoot[C]{\small\thepage}
}

% ---------- Section styling (clean, sharp, no color) ----------
\titleformat{\section}
  {\normalfont\Large\bfseries}{\thesection}{0.8em}{}
\titleformat{\subsection}
  {\normalfont\large\bfseries}{}{0em}{}
\titlespacing*{\subsection}{0pt}{1.4em}{0.6em}

% ---------- Code listings: plain monospace, no color ----------
\lstset{
  basicstyle=\ttfamily\small,
  breaklines=true,
  breakatwhitespace=false,
  columns=fullflexible,
  frame=single,
  framerule=0.4pt,
  rulecolor=\color{black},
  backgroundcolor=\color{white},
  showstringspaces=false,
  xleftmargin=1pt,
  xrightmargin=1pt,
  tabsize=2,
  aboveskip=0.8em,
  belowskip=0.8em
}

% ---------- Helper macros for entries ----------
\newcommand{\entrymeta}[3]{%
  \noindent\textbf{Commit:} \ttfamily#1\normalfont \\
  \textbf{Touches:} \ttfamily\seqsplit{#2}\normalfont \\
  \textbf{Status:} #3
  \par\vspace{0.5em}
}

\newcommand{\entrysep}{\vspace{0.4em}\noindent\rule{\linewidth}{0.3pt}\vspace{1.1em}}

% A printed name with a styled, greyed-out "signature" underneath.
\newcommand{\devcredit}[2]{%
  \begin{minipage}[t]{3.2cm}
    \centering
    {\small #1}\\[0.35em]
    {\color{gray!65}\fontfamily{pzc}\selectfont\large #2}
  \end{minipage}%
}

\setlist[itemize]{leftmargin=1.4em, itemsep=0.2em}
\setlist[enumerate]{leftmargin=1.6em, itemsep=0.2em}

\begin{document}

% ================= TITLE PAGE =================
\begin{titlepage}
\thispagestyle{empty}
\centering
\vspace*{\fill}

{\small\textls[280]{KERNEL BORDERLANDS}}

\vspace{1.6cm}

{\fontsize{36}{44}\selectfont\bfseries Session Report}

\vspace{0.8em}

{\normalsize Attack-Lab Hardening, Feature Schema v2, and the KB-AL-150 CVE Registry}

\vspace{1.2cm}

{\small Compiled for October 31 -- November 1, 2026 (20:12 -- 00:57, one continuous session)}

\vspace{1.2cm}

\begin{minipage}{0.82\linewidth}
\centering
\small
\textbf{Scope} \\[0.3em]
Eight commits across the attack-lab, the Go control plane, and the project's
documentation: two new network-using benign scenarios and a third
90-execution corpus; the \texttt{dest\_class}/\texttt{dest\_port} feature
schema v2 change; the \texttt{Attack-LabTest-Suite.md} benchmark-framework
proposal; and the KB-AL-150 CVE registry, built to 148 of 150 entries over
five research passes. Every figure in this report is pulled directly from
the commits themselves (\texttt{git show}, \texttt{git diff --stat}) and
from the files as committed, not reconstructed from memory.
\end{minipage}

\vspace{1.4cm}

{\small \scshape Kernel Borderlands Team}\\[0.9em]
\devcredit{Karthik}{2311CS040080}\hspace{0.6cm}\devcredit{PardhuVarma}{2311CS040089}\hspace{0.6cm}\devcredit{Tejaswini}{2311CS040077}\hspace{0.6cm}\devcredit{Rupa}{2311CS040082}

\vspace{1.5cm}
\end{titlepage}

% ================= TOC =================
\tableofcontents
\thispagestyle{plain}
\clearpage

% ================= INTRODUCTION =================
\section*{Introduction}
\addcontentsline{toc}{section}{Introduction}

This report covers one continuous working session, starting the evening of
October~31, 2026 and running into the early hours of November~1. Eight
commits landed on \texttt{main} and were pushed to
\texttt{github.com/pardhuvarmax/kernel-borderlands}; none carry Claude
attribution, per standing project policy. The work falls into three
threads, covered in order below: closing a gap in the attack-lab's matched
benign/attack scenario pairs (Section~\ref{sec:attacklab}); a feature
schema change that adds destination-class telemetry to the learned-policy
feature contract without touching the C/Go wire contract
(Section~\ref{sec:schema}); and the KB-AL-150 CVE registry, a
148-of-150-entry metadata catalog built as the first step toward a
CVE-based detection benchmark (Section~\ref{sec:cve}). A fourth section
records what remains open. No exploit code was written at any point in
this session; every CVE-registry entry is explicitly unvalidated
(``category D'') metadata only.

\clearpage

% ================= TIMELINE =================
\section{Timeline}
\label{sec:timeline}

\begin{longtable}{@{}>{\ttfamily}p{1.7cm} p{2.1cm} p{7.6cm} p{2.6cm}@{}}
\toprule
\textbf{Commit} & \textbf{Time} & \textbf{Message} & \textbf{Diffstat} \\
\midrule
\endhead
8254628 & 20:12 & \normalfont attack lab: benign scenarios that read secrets AND use the network & 4 files, +115 \\
1e13978 & 21:45 & \normalfont datasets: third attack-lab corpus (90 executions) with network-using benign pairs & 13 files, +780/-6 \\
283767f & 22:17 & \normalfont feature schema v2: dest\_class/dest\_port from existing NetFlow, no wire change & 8 files, +445/-9 \\
5dc8da4 & 22:56 & \normalfont attack lab: KB-AL-150 CVE registry, cycle 1 of 4 (37 entries) & 3 files, +750 \\
2da05a0 & 23:35 & \normalfont attack lab: KB-AL-150 CVE registry, cycle 2 of 4 (25 entries, total 62/150) & 3 files, +551/-47 \\
4ca13d1 & 00:02 & \normalfont attack lab: KB-AL-150 CVE registry, cycle 3 of 4 (28 entries, total 90/150) & 3 files, +543/-38 \\
678fdae & 00:48 & \normalfont attack lab: KB-AL-150 CVE registry, cycle 4 of 4 (21 entries, total 111/150) & 3 files, +439/-51 \\
68e9920 & 00:57 & \normalfont attack lab: KB-AL-150 CVE registry, fifth research pass (37 entries, total 148/150) & 3 files, +672/-63 \\
\bottomrule
\end{longtable}

\noindent\textbf{Session total:} 36 files changed, 4{,}950 insertions, 11
deletions (\texttt{git diff --stat 75a5330..68e9920}). All eight commits
are authored as \texttt{PardhuVarma Konduru}
\seqsplit{\texttt{<140948604+pardhuvarmax@users.noreply.github.com>}}
and pushed to \texttt{main}.

\clearpage

% ================= SECTION 1: ATTACK LAB =================
\section{Attack-Lab: Closing the Network-Using Benign Gap}
\label{sec:attacklab}

\subsection*{The problem}
Two earlier matched-pair scenarios (\texttt{exfil\_staging} vs.
\texttt{backup\_etc}, \texttt{credential\_dump} vs.
\texttt{credential\_audit}) were designed so the benign half used the same
tool as its attack twin. But neither benign half made a network
connection, so ``has network activity'' separated the pairs by
construction -- not because of anything KB actually detected.

\entrymeta{8254628}{scripts/attack-lab/payloads.py, README.md}{Merged}
\textbf{Fix.} Added a loopback \texttt{\_Sink} listener and \texttt{\_upload()}
helper, then two new scenarios: \texttt{benign\_backup\_remote} (the same
\texttt{sudo tar czf /etc}, then a real upload to a loopback ``backup
host'') and \texttt{benign\_credential\_report} (a counts-only
\texttt{/etc/shadow} audit whose summary is sent to a loopback monitoring
endpoint). A new test
(\texttt{test\_network\_using\_benign\_pair\_members\_read\_secrets\_and\_}
\texttt{connect\_like\_the\_attacks}) brought the lab suite to 14 passing
tests before the capture; the full scenario set grew to 18 (90 executions
at 5 repeats).

\entrysep

\entrymeta{1e13978}{datasets/attack-lab/20261031T201927/*, scripts/attack-lab/tree\_oracle.py}{Merged}
\textbf{Capture.} 90 executions, 321{,}662 events, zero stream gaps, seed
404, \texttt{kbd} observe-only (no automatic containment), sensor
\texttt{exit-state=all}. Run on the dev VM
(\texttt{UbuntuVM-eBPF}, kernel 6.8.0-139-generic). The capture's drop-in
systemd overrides were removed automatically on completion.

\begin{table}[h]
\centering
\small
\begin{tabular}{@{}l l c c c@{}}
\toprule
\textbf{Family} & \textbf{Kind} & \textbf{n} & \textbf{Tree: file \& network} & \textbf{Single process: both} \\
\midrule
exfil\_staging      & attack & 5 & 5/5 & 0/5 \\
backup\_etc         & benign & 5 & 0/5 & 0/5 \\
backup\_remote      & benign & 5 & \textbf{5/5} & 0/5 \\
\addlinespace
credential\_dump    & attack & 5 & 5/5 & 0/5 \\
credential\_audit   & benign & 5 & 0/5 & 0/5 \\
credential\_report  & benign & 5 & \textbf{5/5} & 0/5 \\
\bottomrule
\end{tabular}
\caption{\texttt{tree\_oracle.py}: whether any single process, or the
whole process tree, shows both file evidence (\texttt{dim\_file} $\geq$ 50)
and network evidence. Grouping uses ground truth (ancestry attribution),
so this is an oracle result, not what a deployed rule could achieve.}
\end{table}

\textbf{Finding.} ``File plus network evidence in one process tree'' is
5/5 for both attacks \emph{and} both new network-using benign scenarios.
The earlier result (section omitted here, see
\texttt{datasets/attack-lab/20261031T193756}) was an artifact of the
benign halves never touching the network -- a tree-aggregation rule of
this shape would false-alarm on a real offsite backup. No rule or
feature-schema change was made in response; that is explicitly left as a
decision for the project maintainers (CLAUDE.md's documentation-precedence
rule: conflicts between findings and existing behaviour are surfaced, not
silently resolved).

\textbf{Fixed-rule and baseline detail.} Across the 90-execution corpus,
the fixed rule engine flags attacks 30/40 and benign runs 25/50 at
SUSPICIOUS or higher; only \texttt{process\_injection} reaches
BORDERLANDS (5/5). On the held-out split (\texttt{evasion} and
\texttt{admin\_privileged} families held out entirely, seed 7): fixed
rules 15/19 attacks, 13/23 benign; logistic regression 15/19, 8/23;
gradient-boosted tree 15/19, 13/23; the tree without zone/score features
13/19, 15/23. Sample sizes are small enough that these are descriptive,
not performance claims -- the corpus's own \texttt{baselines\_report.txt}
says so explicitly, and this report repeats that caveat rather than
dropping it.

\clearpage

% ================= SECTION 2: SCHEMA V2 =================
\section{Feature Schema v2: Destination Class Without a Wire Change}
\label{sec:schema}

\entrymeta{283767f}{kb-control-plane/internal/controlplane/\{controlplane.go, controlplane\_test.go\}, kb-aads/features/schema.py, kb-aads/tests/test\_feature\_schema.py, docs/architecture/kbd-contracts.md, docs/development/control-aads/learned-policy-retraining.md}{Merged}

\subsection*{Why this was asked for}
Directly motivated by the Section~\ref{sec:attacklab} finding: if
``file plus network'' can't separate attacks from legitimate backups, can
the \emph{destination} the network traffic goes to help? The operator
approved adding \texttt{dest\_class} and \texttt{bytes\_sent} as new
sensor features, with the explicit caveat that a wire-contract change (C
sensor $\leftrightarrow$ Go control plane) needs both sides to move
together.

\subsection*{What shipped: \texttt{dest\_class} / \texttt{dest\_port}}
\texttt{NetFlowMsg} (wire \texttt{msg\_type=9}) already carries
\texttt{daddr}/\texttt{dport} for every \texttt{connect()} the sensor
observes -- it just wasn't joined into the learned-policy feature export.
\texttt{OnProcessState} now caches the most recent \texttt{NetFlow} sample
per PID (\texttt{netFlowCache}, a \texttt{sync.Map}, evicted on process
exit alongside the existing \texttt{commCache}) and attaches two new
fields to the exported metadata:

\begin{lstlisting}[language=Go]
// classifyDestination buckets an IPv4 destination into a coarse class:
// 0=none 1=loopback 2=private 3=link-local 4=multicast 5=public.
// NetFlowMsg.Daddr keeps the original network-byte-order IP bytes, but
// wire.go's parseNetFlow reads them with binary.LittleEndian.Uint32, so
// the Go value has those bytes reversed from a "normal" IPv4 numeric
// value -- LittleEndian.PutUint32 undoes that.
func classifyDestination(daddr uint32) int {
    b := make(net.IP, 4)
    binary.LittleEndian.PutUint32(b, daddr)
    ...
}
\end{lstlisting}

\texttt{feature\_schema} moved from \texttt{"1"} to \texttt{"2"}; the
existing checkpoint gate (\texttt{gate\_checkpoint} in
\texttt{kb-aads/features/schema.py}) automatically refuses any
v1-trained checkpoint unless \texttt{KB\_ALLOW\_LEGACY\_POLICY=1} is set,
so nothing old is silently fed the new inputs.

\subsection*{What did \emph{not} ship: \texttt{bytes\_sent}}
\texttt{connect()} carries no payload-size field. Real byte accounting
needs new sensor-side instrumentation on \texttt{write()}/
\texttt{sendto()}/\texttt{sendmsg()}, correlating file descriptors to the
socket already tracked for \texttt{NetFlow}, a real wire-contract bump
(\texttt{KB\_WIRE\_VERSION} past 3, a new \texttt{msg\_type}), and
batched delivery from day one -- the project's own journal-flood incident
(CHANGELOG, 1.1.0) and the Go-TLS uprobe bug that once crashed
\texttt{tailscaled} and froze the lab VM are the reasons this was scoped,
not built, in one pass. The scoped plan is written up in
\texttt{docs/development/control-aads/learned-policy-retraining.md}~\S3b
and \texttt{docs/architecture/kbd-contracts.md}, flagged as a separate,
C+Go-coordinated task.

\subsection*{Verification}
\texttt{go test ./...} passes across the whole \texttt{kb-control-plane}
module (9 packages), including two new tests
(\seqsplit{\texttt{TestOnProcessState\_EmitsFeatureSchemaV2}},
\seqsplit{\texttt{TestOnProcessState\_DestClassFromPriorNetFlow}}); the Python
schema test file passes 6/6 including a new
\texttt{dest\_class}/\texttt{dest\_port} round-trip test. Both were
re-run on the lab VM (not just locally) before this report was written.

\subsection*{Known limit, stated plainly}
\texttt{dest\_class} is coarse by construction -- it tells you the
\emph{kind} of address (loopback/private/public/\ldots), not its
\emph{identity}. It cannot by itself distinguish a legitimate backup host
from an attacker's host at the same kind of address. A real fix needs a
destination-identity allowlist layered on top, which is a policy
decision for the maintainers, not a sensor change.

\clearpage

% ================= SECTION 3: CVE REGISTRY =================
\section{KB-AL-150 CVE Registry}
\label{sec:cve}

\subsection*{Origin}
Following an operator-provided design (three validation stages --
observe-only, containment-only, active defense; two execution pathways --
manual and AADS; a 150-CVE registry as the shared evaluation corpus),
written up as \texttt{docs/development/Attack-LabTest-Suite.md} earlier
in this session's parent conversation. This section covers only the
registry itself: a metadata catalog, built over five research passes,
that will eventually let each stage/pathway combination be evaluated
against the same set of real vulnerabilities.

\subsection*{Ground rules, held for all five passes}
\begin{itemize}
\item \textbf{No exploit code.} \texttt{registry.yaml} is metadata only --
  CVE ID, component, class, severity, a one-paragraph summary, affected
  versions, fix version, references, and a provenance tag. Nothing
  executable.
\item \textbf{Every entry is category D} (``metadata-only reference, not
  yet experimentally validated'' per the Attack-Lab doc's \S5) -- none of
  these 148 CVEs have been reproduced, simulated, or replayed.
\item \textbf{Sourced or dropped.} Every entry is either confirmed via web
  search in-session (\texttt{provenance:}
  \seqsplit{\texttt{verified-this-session}}, with a source link or
  citation) or a long-established, heavily documented public-record CVE
  (\texttt{provenance:} \seqsplit{\texttt{established-public-record}}).
  Candidates found without a clean, attributable CVE number were not
  added.
\item \textbf{\texttt{expected\_observations} and \texttt{benign\_twin}
  are \texttt{pending} for all 148 entries.} Filling them in needs each
  CVE's technical analysis read in full and a benign twin designed
  around the same subsystem -- the lesson from
  Section~\ref{sec:attacklab} applied prospectively. That is real
  per-entry work, explicitly deferred.
\item \textbf{\texttt{lab.reproduction\_available} is \texttt{false}
  for all 148.} No disposable guest on an affected kernel exists yet.
\end{itemize}

\subsection*{Per-cycle, per-class breakdown}

\begin{table}[h]
\centering
\scriptsize
\begin{tabular}{@{}l r r r r r r r r@{}}
\toprule
\textbf{Class} & \textbf{Target} & \textbf{C1} & \textbf{C2} & \textbf{C3} & \textbf{C4} & \textbf{C5} & \textbf{Total} & \textbf{vs.\ target} \\
\midrule
privilege-escalation               & 20 & 5 & 1 & 2 & 5 & 5 & 18 & $-2$ \\
memory-corruption                  & 20 & 4 & 5 & 2 & 5 & 4 & 20 & $0$ \\
use-after-free                     & 15 & 5 & 5 & 0 & 1 & 4 & 15 & $0$ \\
race-conditions                    & 15 & 4 & 2 & 2 & 2 & 5 & 15 & $0$ \\
filesystem                         & 15 & 4 & 1 & 5 & 3 & 3 & 16 & $+1$ \\
network-stack-netfilter            & 15 & 5 & 2 & 0 & 3 & 4 & 14 & $-1$ \\
container-isolation-escape         & 15 & 5 & 1 & 5 & 0 & 5 & 16 & $+1$ \\
syscall-ioctl-kernel-interface     & 15 & 0 & 2 & 10& 0 & 2 & 14 & $-1$ \\
information-disclosure             & 10 & 3 & 2 & 1 & 1 & 3 & 10 & $0$ \\
credential-permission-manipulation & 10 & 2 & 4 & 1 & 1 & 2 & 10 & $0$ \\
\midrule
\textbf{Total} & \textbf{150} & \textbf{37} & \textbf{25} & \textbf{28} & \textbf{21} & \textbf{37} & \textbf{148} & \textbf{$-2$} \\
\bottomrule
\end{tabular}
\caption{Entries per research pass (C1--C5) by vulnerability class. Six of
ten classes hit their target exactly; two are one over (harmless -- 150
per class was a target, never a hard cap); three are short by one or two
each, left that way rather than filled with weak sources.}
\end{table}

\textbf{Provenance split:} 126 entries \texttt{verified-this-session}
(confirmed via web search, source cited), 22 \texttt{established-public-record}
(long-documented classics -- Baron Samedit, PwnKit, Dirty~COW, Dirty~Pipe,
the netfilter \texttt{x\_tables} OOB write, the \texttt{route4} UAF, the
AF\_PACKET UDP-fragmentation UAF, the CAN~BCM race, several legacy
futex/DCCP/Podman flaws, among others -- not re-searched given how
thoroughly they are already documented elsewhere).

\subsection*{Representative entries}
A small sample across the five passes, illustrating the range of
subsystems covered (full detail for all 148 is in
\texttt{registry.yaml}):

\begin{table}[h]
\centering
\scriptsize
\begin{tabular}{@{}l p{3.1cm} p{7.3cm}@{}}
\toprule
\textbf{CVE} & \textbf{Subsystem} & \textbf{One-line summary} \\
\midrule
CVE-2022-0847 & pipe / page cache & Dirty Pipe: missing \texttt{pipe\_buffer} flag init lets a local user overwrite read-only files. \\
CVE-2021-4034 & polkit \texttt{pkexec} & PwnKit: empty \texttt{argv[0]} falls through to attacker-controlled env/PATH. \\
CVE-2026-23111 & netfilter \texttt{nf\_tables} & A single removed \texttt{!} inverted abort-path activation logic, causing a UAF; public exploit, $>$99\% reliability claimed by one source. \\
CVE-2024-21626 & \texttt{runc} & ``process.cwd and leaked fds'' container breakout (``Leaky Vessels''). \\
CVE-2025-31133/52565/52881 & \texttt{runc} & Three chainable flaws (Aleksa Sarai) that together grant admin-level host access via \texttt{RUN~--mount}. \\
CVE-2023-0386 & OverlayFS & Unmapped uid/gid in \texttt{ovl\_copy\_up\_one} lets a copied setuid binary keep file capabilities. \\
CVE-2026-31431 & AF\_ALG crypto API & ``Copy Fail'': page-cache corruption of setuid binaries; CISA-acknowledged active exploitation. \\
CVE-2023-2163 & eBPF verifier & Incorrect branch pruning marks unsafe paths safe; arbitrary kernel r/w, container escape. \\
\bottomrule
\end{tabular}
\caption{A representative slice of the 148-entry registry.}
\end{table}

\subsection*{Entries carrying extra sourcing caveats}
Eight entries, flagged in-file, need a second look before any downstream
use:

\begin{itemize}
\item \textbf{CVE-2026-31635} (``DirtyDecrypt'', C3) -- a single,
  unconfirmed source claimed no vendor patch existed at time of search.
\item \textbf{CVE-2026-96812} (gVisor, C3) -- NVD listed it as
  ``Received'' (not yet fully analyzed) at time of search.
\item \textbf{CVE-2022-3567} (C4) -- thinly sourced; treat as a lead, not
  a confirmed technical detail.
\item \textbf{CVE-2024-26594} (C4) -- the source found didn't name the
  affected subsystem; component/subsystem are unconfirmed.
\item \textbf{CVE-2026-43198} (C4) -- severity disputed across sources
  (9.8 vs.\ 7.0).
\item \textbf{CVE-2026-43499} (``GhostLock'', C5) -- a third-party
  blog's ``97\% stable exploit'' claim is unverified; a later search in
  the same pass found the public PoC is Android-specific and causes a
  crash, not demonstrated privilege escalation on mainline Linux.
\item \textbf{CVE-2026-90203} (squashfs, C5) -- disclosure date and CVSS
  not independently confirmed.
\item \textbf{CVE-2026-43502} (``ZcopyReaper'', C5) -- exact fixing
  commit not confirmed.
\end{itemize}

\noindent Separately, three entries are \textbf{remote-reachable} rather
than local-only, unlike most of the registry: \textbf{CVE-2023-52440} and
\textbf{CVE-2026-31705} (both \texttt{ksmbd}), and \textbf{CVE-2026-43198}
(TCP/IPv6). The attack-lab's isolation rules apply with extra weight to
any future reproduction work on these three.

\subsection*{Why 148, not 150}
Every one of the final two slots found in the fifth pass was either a
duplicate of an entry already in the registry or came with sourcing too
thin to meet the bar held for the other 148. Padding to an exact round
number with an unverified CVE ID would have been strictly worse than
stopping two short: a wrong entry here doesn't just sit idle, it
propagates into every benchmark run and paper citation built on top of
it. The two-short state is recorded, not hidden, in both this report and
\texttt{registry.yaml}'s own README.

\clearpage

% ================= SECTION 4: OPEN ITEMS =================
\section{Open Items, In Priority Order}
\label{sec:open}

\begin{enumerate}
\item \textbf{Reproduction scoping.} Pick 3--5 registry entries with a
  clear single-mechanism path and an existing public PoC (the
  \texttt{nf\_tables} UAF chain, Dirty~Pipe, Dirty~COW, PwnKit, Baron
  Samedit are the strongest candidates) and scope what a disposable guest
  needs for each: kernel version, image source, snapshot-and-discard,
  no network beyond the harness. Not started -- this is exploit-adjacent
  work requiring the operator's go-ahead per CVE.
\item \textbf{\texttt{expected\_observations} / \texttt{benign\_twin}}
  for those first few entries, designed before any reproduction code is
  written.
\item \textbf{Destination-identity allowlisting}, layered on top of
  \texttt{dest\_class}, to actually separate legitimate backups from
  exfiltration at the same address class -- a policy decision for the
  maintainers.
\item \textbf{\texttt{bytes\_sent}} -- scoped (Section~\ref{sec:schema}),
  not started; needs new eBPF instrumentation and a coordinated wire
  bump.
\item \textbf{Closing the last 2--4 CVE-registry entries}
  (privilege-escalation $-2$, network-stack-netfilter $-1$,
  syscall-ioctl-kernel-interface $-1$) -- low priority next to items 1--2.
\item \textbf{Tree-level evidence aggregation in the fixed rule engine} --
  still an open spec-level decision; per-process evidence is not joined
  across a process tree today, and changing that touches the state
  machine.
\end{enumerate}

\vspace{1em}
\noindent None of these were started after this report was requested;
they are recorded exactly as they stood when the session paused.

\end{document}
